\section{Methodologies}



\begin{table*}[t]
%       \rotatebox{270}{
\centering
\begin{tabular}{lccccc}
        Name  &   Wait & \#Txn-Types & \#WH & \#Workers & Cross-WH (\%) \\
    \toprule
        Vanilla~\cite{tpccurl} & Yes & 5 & Configurable & 10 per WH & 9.5\% and 15\%   \\
    \midrule
        Calvin~\cite{Thomson2012Calvin,ren2014evaluation,calvin-repo}   &  No  &  2   & 20 per server & 4 per server   & 0-100\%   \\
        Silo~\cite{Tu2013Silo,siloSRC}       &  No  &   5  & 1-32 & 1 per WH &  0-100\%   \\
        DrTM~\cite{Wei2015FIT, DrTMSRC}           &  No    &   5   &  8 per Server  &  1 per WH   &  0-100\%          \\
        Janus~\cite{Mu2016Consolidating,janusSRC}    & No  &  5 & 6 & 1-10K   &  9.5\% and 15\%        \\
        Cicada~\cite{Lim2017Cicada,cicadaSRC}     & No &  5    &  1-28  & 1-28 &   9.5\% and 15\%       \\
        GAM~\cite{Cai2018GAM, GamSRC}                &   No   &  2    &  4 per server  & 1 per WH    &  0-100\%          \\
        Star~\cite{Lu2019Star, StarSRC}       &  No  &  2 & 12 per server & 1 per WH & 0-100\%   \\
        Aria~\cite{Lu2020Aria,ariaSRC}        & No    &  2    &  1-180 &   1 per WH   & 0-100\%   \\
    \bottomrule
\end{tabular}
\caption{TPC-C settings used by different systems.}
\label{tab:tpcc}
\end{table*}

 \begin{table*}[t]
 %       \rotatebox{270}{
 \centering
\begin{tabular}{lcm{3cm}cc}
        Name     & \#KV pairs & KV size & R/W ratio & Distribution  \\
    \toprule
        Vanilla~\cite{ycsburl}  & 120M & 1KB & 50:50 or 95:5 or 100:0 & Uniform/Zipfian (0.99)  \\
     %   Facebook Study~\cite{Atikoglu2012WAL}  & ?? & Key 20B-40B (median); Value 50B-500B (median) & ?? & ??  \\
    %    Twitter Study~\cite{Yang2020Twitter}  & ?? & A few bytes to tens of KB & 9:1 (average) & Zipfian 0.9 and 1.4   \\
    \midrule
        Star~\cite{Lu2019Star, StarSRC}      &   2.4M & 8B/100B  &  90:10   &  Uniform   \\
        Silo~\cite{Tu2013Silo,siloSRC}       & 160M  & 8B/100B  &  80:20  &  Uniform \\
        Cicada~\cite{Lim2017Cicada,cicadaSRC}      & 10M  & 8B/100B  & 95:5/50:50 & Zipfian  \\
        HERD~\cite{Kalia2014HERD, HERDSRC}         & 480M  & 16B/(4B - 1KB)  &  50:50   & Uniform/Zipfian    \\
        TAPIR~\cite{Zhang15Building,tapir-repo}     & 1M & 8B/1KB & 50:50 & Zipfian  \\
    \bottomrule
\end{tabular}
\caption{YCSB settings used by different systems.}
\vspace{-.1in}
\label{tab:ycsb}
%}
\end{table*}

We focus on TPC-C and YCSB, two of the most popular benchmarks used in our
community, because there is a general consensus about what parameters to
tune for these two. Then among the works that are evaluated under these two benchmarks,
we selected 10 that we can successfully reproduce, as listed
in Table~\ref{tab:tpcc} and Table~\ref{tab:ycsb} (note that some systems are evaluated under
both). Since our hardware is not exactly the same as those used by these works, we consider ``success''
as our reproduced numbers are reasonably close to those in the original paper and the conclusion of the
paper still holds. For these 10 works, we are able to reproduce most of their results reported in the corresponding papers.
In this work, we mainly focus on the throughput of these systems and we plan to
explore latency in the future.
%\todo{Among them, Cicada has got all the badges of artifact evaluation; the others
%do not partially because some were published before the movement of artifact evaluation, but
%since we are able to reproduce most of their results, we believe they could pass today's artifact
%evaluation.}

Note that this list is by no means a complete list of all reproducible works: we have to skip ones that
require special hardware; we have not got a chance to try most
recent publications, which we plan to explore in the future.

For a very brief introduction, among these 10 works, Silo and Cicada are single-node in-memory
transaction processing engines; Calvin, DrTM, GAM, Star, and Aria support distributed transactions;
TAPIR and Janus further support geo-distributed transactions; HERD is a distributed
key-value store. DrTM, GAM, and HERD use RDMA techniques to improve performance.


%Note that some of these systems are published before the movement of artifact evaluation,
%but since we are able to reproduce their results, we believe they could pass today's artifact evaluation.
%Second, we do not try other benchmarks because, compared to TPC-C and YCSB,
%there is less consensus about which parameters should be tuned for these benchmarks,
%but we plan to explore more benchmarks and systems in the future.



After reproducing the original results of these systems, we try to test them under a broader range of settings.

\vspace{-.1in}
\paragraph{Parameters of TPC-C.}
TPC-C~\cite{tpccurl} simulates the database
of a wholesale company, which includes a number of warehouses.
TPC-C consists of five types of transactions: the \emph{New-Order} transaction simulates
the procedure of entering an order. Since \emph{New-Order} transactions targeting the same warehouse
need to read and update the same \emph{next-order-ID}
value, they are the major contention point of TPC-C. The \emph{Payment} transaction 
updates the customer's balance and reflects the payment on the district and warehouse sales statistics.
The other three types are less frequent so we skip their details.

Table~\ref{tab:tpcc} shows the common parameters tuned by different works: ``Wait'' identifies
whether a customer will wait a while before starting the next transaction; ``Txn-Types'' identifiers
how many types of transactions to run (5 for all types and 2 for \emph{New-Order} and \emph{Payment});
``\#WH'' identifies how many warehouses to use (also called ``scale factor'' in TPC-C);
``\#Workers'' identifies how many concurrent customers may access a warehouse; and finally both
\emph{New-Order} and \emph{Payment} have a chance to access a remote warehouse, and such chance
is identified by ``Cross-WH''\footnote{TPC-C uses different chances for \emph{New-Order} and \emph{Payment},
but since many prototypes tune them together, we only list one here.}.

Vanilla TPC-C only allows to tune the number of warehouses:
with a long wait time per customer and a restriction of 10 customers per warehouse, vanilla
TPC-C puts a very low limit on the throughput one can get per warehouse. Therefore, to get a high
throughput, one needs to use many warehouses, which results in a low-contention and I/O heavy workload.
Since many research prototypes focus on addressing contentions, they remove wait time and use a small number
of warehouses to reach a high contention level.
It is out of the scope of this paper to determine whether such tuning represents any realistic workload,
since we cannot find any studies about transaction workload and it is unclear whether the 30-year old
TPC-C can still represent a realistic workload today. As a result,
an extensive evaluation is perhaps a safe choice.
%It is unclear whether the 30-year old warehouse-based TPC-C can still represent any realistic workloads
%today. Even so, from Table~\ref{tab:tpcc}, we can see most research prototypes deviate from the original
%settings of TPC-C and it's well-known that such deviation actually significantly changes the behavior of
%TPC-C: 

For an extensive evaluation, we 1) do not use waiting since no systems use it; 2) try both 2 and 5
transaction types if the target system allows; 3) tune the number of warehouses from 1 to a large value
allowed by DRAM (all target systems assume all data can fit into DRAM); 4) tune the number of workers
from 1 to a large value; 5) tune cross-warehouse rate from 0\% to 100\%.
We list detailed parameter values in Table~\ref{tab:summary}, but note that some systems disallow
the tuning of certain parameters.

\vspace{-.1in}
\paragraph{Parameters of YCSB.}

Yahoo! Cloud Serving Benchmark (YCSB) is a benchmark primarily targeting
 NoSQL database systems with a key-value like interface~\cite{Cooper2010YCSB}.
 The commonly tuned parameters include number of key-value pairs, size of a key-value pair,
 read/write ratio, and the distribution of keys.
 
Table~\ref{tab:ycsb} lists the parameter values used by the vanilla YCSB paper
and those used by research prototypes. Furthermore, we find a number of
studies about production key-value stores~\cite{Atikoglu2012WAL,Yang2020Twitter,Nishtala2013Scaling,Berg2020CacheLib},
and these studies show that industrial companies can have a variety of
workloads. To quote a few conclusions from the recent Twitter study~\cite{Yang2020Twitter}:

\vspace{.05in}
\noindent
---``in addition to the well-known use case of serving read-heavy workloads, a substantial
number of Twemcache clusters are used to serve write-heavy workloads.''
%``most of the existing systems, optimizations and research assume a
%read-heavy workload''; ``This calls for future research on designing systems
%and solutions that consider performance on write-heavy workloads''.

\vspace{.05in}
\noindent
---``a recent work from Facebook~\cite{Berg2020CacheLib} suggested that in-memory caching workloads
may not follow Zipfian distribution.''; ``Measuring all Twemcache workloads, we observe the majority
of the cache workloads still follow Zipfian distribution.''; ``most of the $\alpha$ values (of Zipfian distribution) are in the range
from 1 to 2.5.''

\vspace{.05in}
\noindent
---``around 85\% of Twemcache clusters have a mean key size smaller than 50 bytes, with
a median smaller than 38 bytes''; ``the mean value size falls in the range from 10 bytes to 10 KB, and
25\% of workloads show value size smaller than 100 bytes,
and median is around 230 bytes.''

The versatility of production workloads, both within a single organization
and across different organizations, is exactly one of the motivations for our pursuit of extensive
evaluation: even if a work is optimized for one important workload, it would
be helpful to know how it performs on other workloads.

For an extensive evaluation, we try to 1) tune the key-value size from 8B to 1KB;
2) tune the number of key-value pairs from 128K to the maximal allowed by DRAM;
3) tune read/write ratio from 100:0 to 0:100; and 4) try both uniform and
Zipfian distribution and try different $\alpha$ values for the Zipfian distribution.

\vspace{-.1in}
\paragraph{System parameters.} For distributed systems, we measure
system scalability by tuning the number of servers. For systems supporting
replication, we also tune the number of replicas. However, due to resource
limitations, we do not try a large number of servers.


\paragraph{Testbed.}
We run our experiments on CloudLab~\cite{cloudlab} and a local cluster.
We try to use hardware that is close to the one used by the original paper.
On CloudLab, we use m510 machines to
run Aria experiments, c6320 machines to run Silo and Cicada experiments, rs630
to run Janus experiments, c220g1 machines to run TAPIR experiments, d430 to run
Calvin experiments, c6220 machines to run Star and GAM experiments, and r320 machines to run HERD experiments.
We test DrTM on a local cluster, in which each node is equipped
with a 28-core Intel Xeon E5-2680 CPU and 128 GB of RAM. The cluster is interconnected by a 100 Gbps Mellanox CX-4 adapter.
We use OpenFabrics Enterprise Distribution (OFED) version 4.9.
